% !TEX root = ../thesis.tex

\chapter{Conclusion}\label{ch:conclusion}

\section{Summary of the Work}

This bachelor thesis addressed the problem of binary semantic segmentation of archaeological features in airborne LiDAR-derived raster data. The work was motivated by the fact that LiDAR data provide detailed information about terrain morphology, but their interpretation in archaeology still depends strongly on expert visual inspection. Since large LiDAR datasets are time-consuming to inspect manually, the thesis investigated whether deep learning can support archaeological interpretation by producing reproducible candidate maps of potential archaeological features.

The main objective was not only to train a segmentation model, but to design, implement, and experimentally evaluate a reproducible workflow for semantic segmentation in two different archaeological and geographical domains. The first domain consisted of compact Maya archaeological structures in Guatemala. The second domain consisted of terrace-like archaeological features in Slovakia. These two domains differ in object morphology, spatial resolution, annotation type, and terrain context. Their joint use therefore made it possible to evaluate not only the performance of the selected models, but also the sensitivity of the whole workflow to changes in data domain.

The thesis first analysed traditional and semi-automatic methods used in archaeological LiDAR interpretation and then reviewed the role of deep learning in remote sensing and archaeology. The analysis showed that deep learning methods are promising for archaeological prospection, but their practical usefulness depends on many methodological decisions. These include the choice of input representation, the conversion of expert annotations into training masks, the sampling strategy, the loss function, the model architecture, threshold optimization, post-processing, and the way in which the resulting maps are evaluated. For this reason, the work focused on building a controlled experimental framework rather than on presenting a single isolated model result.

A complete processing workflow was created. It included preparation of LiDAR-derived input representations, conversion of reference annotations into hard and soft masks, construction of tiled datasets, definition of train-validation strategies, model training, threshold-based inference, quantitative evaluation, and qualitative GIS-based inspection. The implemented workflow made it possible to compare several controlled variables under consistent conditions. In the Maya domain, the compared input representations included DEM, RGB-like visualization, RGDEM, and SPS. In the Slovak domain, DEM, ISO, and ANISO representations were evaluated. The work also compared hard and soft label variants, different tile sizes, selected loss functions, and DeepLab-based segmentation architectures.

The best-performing configuration for the Maya dataset was based on DeepLab~V3+ with an ImageNet-pretrained ResNet-101 backbone, SPS input representation, soft labels with a radius of 2~m, Tanimoto with complement loss, and $128 \times 128$ pixel tiles with a stride of 32 pixels. This configuration achieved stronger results than the earlier DeepLab~V3 baseline across the retained pixel-wise metrics. In particular, the foreground IoU increased from 0.24690 to 0.37924, the F1 score from 0.39602 to 0.54992, and the Matthews correlation coefficient from 0.39634 to 0.56307. From the object-level perspective, the selected DeepLab~V3+ configuration detected 2253 out of 4084 reference objects, while the compared DeepLab~V3 baseline detected 1746 objects. The improvement should therefore be interpreted primarily as an increase in detection completeness.

The best-performing configuration for the Slovak terrace dataset was based on DeepLab~V3+ with a ResNet-101 backbone, ANISO input representation, soft labels with a radius of 2~m, BCE+Dice loss, and $512 \times 512$ pixel tiles with a stride of 128 pixels. This configuration achieved the strongest overall performance among the evaluated Slovak variants. However, the Slovak results were substantially more difficult to interpret than the Maya results. The terrace features are narrow, elongated, partially discontinuous, and represented by buffered polylines rather than precise polygonal outlines. As a result, the model often detected broader terrace-like zones rather than individual terrace lines. The obtained foreground IoU of 0.06560 and F1 score of 0.12313 show that the output is not suitable as a final segmentation map. Nevertheless, the object-level hit rate of 16 detected objects out of 24 indicates that the model can still be useful as a candidate localization layer.

The ablation study showed that LiDAR-derived input representation has a clear influence on segmentation quality. On the Maya dataset, SPS achieved the strongest results among the compared representations, while DEM produced the weakest result. This confirms that direct elevation values alone are less informative for this task than representations that emphasize local terrain morphology. The comparison of hard and soft labels showed a more nuanced result. Soft masks did not improve all configurations uniformly, but they often increased recall and reduced precision. This behaviour is useful when the main goal is to avoid missing potential archaeological candidates, but it also increases the risk of false-positive predictions. The tile-size experiments showed that smaller tiles were more suitable for compact Maya structures, while a larger tile size was more appropriate for Slovak terraces, where broader terrain context is needed. The loss-function comparison showed that BCE+Dice produced the strongest metric-oriented result in one comparison group, while Tanimoto with complement provided more stable training behaviour and was used in the strongest Maya configuration. The architecture comparison showed that DeepLab~V3+ was preferable to DeepLab~V3 because it provided a slightly better precision-recall balance and more stable behaviour.

The qualitative error analysis showed that the trained models do not fail randomly. False-positive predictions usually occurred on natural or disturbed terrain forms that visually resembled archaeological structures. False-negative predictions were especially common for small, isolated, weakly preserved, or low-contrast objects. Boundary errors appeared when the model localized a relevant structure but did not reproduce the full reference outline. This behaviour reflects an important property of archaeological LiDAR segmentation: expert annotations do not always correspond exactly to the currently visible terrain expression. The model tends to follow the raster-visible relief signal, while the annotation may represent an interpreted archaeological object. Therefore, some metric penalties reflect not only model error, but also uncertainty in the relation between the terrain representation and the reference geometry.

The archaeological usefulness of the results depends on the intended use case. For the Maya dataset, the model is useful mainly as a candidate detector, especially for medium and large structures. Large structures were detected most reliably, with 400 detected objects out of 472 and a hit rate of 0.8475. Medium-sized structures were detected with a hit rate of 0.6338, while small structures reached only 0.3287. This confirms that object size and terrain expression strongly influence detectability. For the Slovak terrace dataset, the model is not suitable for automatic extraction of precise terrace geometries, but it can highlight probable terrace zones. In both cases, the output should be interpreted as a decision-support layer for expert-guided inspection rather than as an automatic archaeological inventory.

\section{Fulfilment of Objectives}

The main objective of the thesis was fulfilled. A reproducible deep learning workflow for binary semantic segmentation of archaeological features in airborne LiDAR-derived raster data was designed, implemented, and experimentally evaluated. The workflow was applied to two different domains: Maya archaeological structures in Guatemala and terrace-like archaeological features in Slovakia. The experiments demonstrated that the same general workflow can be used in both domains, but also showed that domain-specific adjustment is necessary.

The first partial objective was to analyse the current state of remote sensing, LiDAR-based archaeological prospection, and deep learning methods used for archaeological object detection and semantic segmentation. This objective was fulfilled in the analytical part of the thesis. The work described the transition from traditional remote sensing and manual LiDAR interpretation toward semi-automatic and deep learning-based methods. It also identified the main methodological problems relevant to this thesis: input representation, label uncertainty, class imbalance, and the difference between numerical segmentation quality and archaeological usefulness.

The second partial objective was to prepare LiDAR-derived input representations suitable for CNN-based semantic segmentation. This objective was fulfilled by preparing and evaluating several raster representations derived from the available terrain data. For the Maya dataset, DEM, RGB-like visualization, RGDEM, and SPS were compared. For the Slovak dataset, DEM, ISO, and ANISO representations were used. The results showed that input representation is one of the most important components of the workflow. In the Maya experiments, SPS provided the strongest performance, while DEM was the weakest representation. This confirms that derived terrain visualizations can provide more useful information for segmentation than direct elevation values alone.

The third partial objective was to prepare reference masks for both study domains and to compare hard binary labels with soft label variants. This objective was fulfilled by converting the available expert annotations into raster masks and by constructing both hard and soft label representations. The experiments showed that soft masks are not universally better than hard masks. Their most visible effect was a shift in the precision-recall balance: they often increased recall while reducing precision. This means that soft labels can be useful when the priority is to increase candidate coverage, but they must be used carefully because broader predictions may also increase false positives.

The fourth partial objective was to construct a reproducible tiled dataset pipeline based on manifests, valid-pixel masks, controlled spatial splits, and consistent train-validation strategies. This objective was fulfilled through the dataset construction workflow. Raster data and masks were divided into tiles and organized into controlled experimental subsets. This made it possible to compare model configurations systematically and to reduce the risk that differences in results were caused by changes in data preparation rather than by the tested methodological component.

The fifth partial objective was to train and compare selected semantic segmentation models on the Maya and Slovak datasets under controlled experimental settings. This objective was fulfilled by training DeepLab-based models on both study domains and comparing their performance using consistent metrics. The work compared DeepLab~V3 and DeepLab~V3+ variants, including pretrained and non-pretrained configurations. The strongest Maya configuration used DeepLab~V3+ with an ImageNet-pretrained ResNet-101 backbone, while the strongest Slovak configuration used DeepLab~V3+ with a ResNet-101 backbone trained from scratch. These results show that the most appropriate model setup depends on the domain and the morphology of the target objects.

The sixth partial objective was to evaluate the influence of input representation, label type, sampling strategy, loss function, model architecture, and post-processing on segmentation performance. This objective was fulfilled through the ablation study and the comparison of representative configurations. The experiments showed that SPS was the strongest Maya input representation, that soft labels mainly affected the recall-precision balance, that tile size must be selected according to object morphology and raster resolution, that loss functions differ not only in final metrics but also in training stability, and that DeepLab~V3+ is preferable to DeepLab~V3 in the evaluated setup. The results therefore identify which workflow components had a clear impact and which effects were more limited or domain-dependent.

The seventh partial objective was to assess the obtained results using both quantitative pixel-wise metrics and qualitative inspection in a GIS environment. This objective was fulfilled. Pixel-wise metrics such as foreground IoU, F1 score, precision, recall, balanced accuracy, and Matthews correlation coefficient were used for controlled comparison. At the same time, qualitative GIS-based inspection was used to interpret false positives, false negatives, boundary errors, and successful predictions. This combined evaluation was necessary because pixel-level metrics alone do not fully describe whether a prediction is archaeologically useful.

The eighth partial objective was to identify representative configurations that either improve segmentation quality or clearly demonstrate limitations of the proposed workflow. This objective was fulfilled in the results and discussion. The best-performing Maya and Slovak configurations were identified, and weaker or more ambiguous configurations were discussed through the ablation study. The analysis also identified the main limitations of the workflow, including reference-data uncertainty, extreme class imbalance, limited cross-domain generalization, differences in object morphology, and the need for expert interpretation of full-area inference results.

The research questions defined in the thesis were also addressed. The experiments showed that LiDAR-derived input representation has a substantial effect on segmentation quality, that visually enhanced raster representations can outperform direct DEM input, and that combining or deriving terrain features can improve model performance when the representation emphasizes relevant micro-relief. The comparison of hard and soft labels showed that label representation influences training behaviour and prediction selectivity. The loss-function experiments showed that different objectives lead to different precision-recall behaviour and training stability. The model comparison showed that DeepLab~V3+ is generally preferable in the tested setting, but that architecture alone is not the only determining factor. Finally, the GIS-based analysis confirmed that the best pixel-wise results do not automatically equal a final archaeological interpretation. The most useful outputs are those that can support expert inspection by highlighting candidate areas.

\section{Main Contributions}

The first contribution of this thesis is the design and implementation of a reproducible workflow for semantic segmentation of archaeological features in LiDAR-derived raster data. The workflow connects geospatial data preparation, reference mask generation, tiled dataset construction, model training, threshold-based inference, metric evaluation, and GIS-based interpretation. This is important because the quality of deep learning results in archaeological LiDAR analysis depends not only on the model architecture, but also on many preceding and following methodological steps.

The second contribution is the controlled comparison of several methodological components under a consistent experimental framework. The thesis did not treat the trained model as a black-box result. Instead, it evaluated the influence of input representation, label representation, tile size, loss function, and model architecture. This makes the results more informative than a single final score, because it shows which parts of the workflow influenced the final output and how these effects differed between the Maya and Slovak domains.

The third contribution is the experimental evaluation of different LiDAR-derived input representations. The results demonstrated that SPS was the most effective representation for Maya structures among the tested variants, while DEM was the weakest. This confirms that representations emphasizing local terrain morphology are more useful for this segmentation task than direct elevation values alone. For the Slovak dataset, the best-performing configuration used ANISO, which indicates that the optimal representation depends on the morphology of the target objects and the landscape context.

The fourth contribution is the evaluation of hard and soft label representations for archaeological segmentation. The thesis showed that soft labels can be useful, but their effect is not uniformly positive. They can increase recall and make the model more inclusive, but this may reduce precision and increase false positives. This result is important because archaeological annotations are often uncertain, schematic, or spatially imprecise. Soft labels can partly address this problem, but they do not remove the need for careful interpretation of prediction boundaries.

The fifth contribution is the evaluation of the practical archaeological usefulness of the predictions, not only their pixel-wise quality. For the Maya dataset, the best model detected 2253 out of 4084 reference objects. The size-based analysis showed that the model was much more reliable for large and medium-sized structures than for small structures. This result provides a more archaeologically meaningful interpretation than pixel-level metrics alone. For the Slovak terrace dataset, the model detected 16 out of 24 terrace objects, but the qualitative analysis showed that the output should be interpreted as candidate localization rather than precise delineation.

The sixth contribution is the identification of the main error types in archaeological LiDAR segmentation. The thesis distinguished false positives, false negatives, and boundary errors and related them to terrain morphology, label uncertainty, object size, and visualization effects. This analysis showed that many false positives appear in areas that are visually similar to archaeological structures, while false negatives are common for small or weakly expressed objects. Boundary errors often reflect the mismatch between expert-interpreted object outlines and the visible terrain signal. These observations are important for understanding how model outputs should be used in practice.

The seventh contribution is the comparison of two different archaeological domains within the same methodological framework. The Maya and Slovak datasets differ substantially in morphology and annotation type. The Maya dataset contains compact archaeological structures, while the Slovak dataset contains narrow and elongated terrace-like features. The results showed that a workflow suitable for one domain cannot be transferred directly to another without adjustment. This supports the conclusion that archaeological deep learning workflows require domain-specific validation.

The final contribution is the formulation of a realistic interpretation of the role of deep learning in archaeological LiDAR prospection. The thesis does not claim that the developed models replace expert interpretation. Instead, it shows that semantic segmentation can be useful as a decision-support tool. The strongest results are suitable for highlighting candidate locations and reducing the amount of terrain that must be inspected manually, especially in the case of medium and large Maya structures and broader Slovak terrace zones. The final archaeological interpretation, however, remains dependent on expert review.

\section{Future Work}

Several directions for future work follow directly from the limitations and results of this thesis.

The first direction is the improvement of reference data. The experiments showed that reference annotations strongly influence both model training and evaluation. Future work should therefore focus on producing more consistent expert annotations, documenting uncertainty in the labels, and distinguishing between visible terrain expression and interpreted archaeological extent. For the Slovak terrace dataset, more precise polygonal annotations or several alternative label representations could help clarify whether the model fails because of weak terrain signal, unsuitable architecture, or mismatch between buffered polylines and actual relief morphology.

The second direction is the extension of the evaluation protocol. Pixel-wise metrics are necessary, but they are not sufficient for archaeological interpretation. Future work should further develop object-level and GIS-based evaluation methods that better reflect archaeological usefulness. In particular, candidate localization, distance from reference features, object size, connected components, and expert review categories may provide more informative assessment than foreground IoU alone. Such evaluation would be especially useful for terrace systems, where a prediction can indicate the correct area even if it does not overlap precisely with the reference line.

The third direction is the development of more domain-specific architectures and training strategies. This thesis evaluated selected DeepLab-based models, but the results suggest that archaeological LiDAR segmentation may benefit from architectural adjustments. For compact Maya structures, models that preserve local boundary detail may improve delineation. For Slovak terraces, models better adapted to narrow, elongated, and discontinuous features may be necessary. Future experiments may therefore investigate architectures with stronger multiscale feature fusion, attention mechanisms, higher-resolution decoders, or methods designed specifically for thin linear objects.

The fourth direction is the comparison of semantic segmentation with object detection and instance segmentation approaches. This thesis was limited to binary semantic segmentation. However, the results show that object-level usefulness is sometimes more important than exact pixel-wise overlap. For Maya structures, instance-level methods could help separate individual structures. For Slovak terraces, object detection or line-aware methods could be more appropriate than dense binary segmentation if the goal is candidate localization rather than exact mask extraction.

The fifth direction is improved handling of class imbalance and false positives. Archaeological features occupy only a small fraction of the raster area, and this creates a strong foreground-background imbalance. Future work should further investigate sampling strategies, loss functions, and post-processing methods that improve recall without producing excessive false positives. This is especially important for full-area inference, where a large number of false-positive candidates can reduce the practical usefulness of the model.

The sixth direction is broader cross-domain validation. The two datasets used in this thesis already demonstrated that results depend strongly on data domain. Future work should therefore test the workflow on additional archaeological contexts, additional terrain types, and additional annotation styles. Such validation would help determine which conclusions are specific to the Maya and Slovak datasets and which methodological findings can be generalized to other LiDAR-based archaeological prospection tasks.

The seventh direction is the integration of the model outputs into a human-in-the-loop archaeological workflow. Since the model outputs are most useful as candidate layers, future work should focus on how predictions can be reviewed, corrected, and reused by experts. A practical workflow could combine probability maps, binary masks, object candidates, uncertainty indicators, and GIS-based editing. Expert corrections could then be incorporated into later training iterations, gradually improving the dataset and the model.

The final direction is the possible extension beyond raster-derived products. This thesis worked only with LiDAR-derived raster representations and did not process the original point clouds directly. Future research may compare raster-based segmentation with methods that use point-cloud information or additional terrain derivatives. Such work would be beyond the scope of the present thesis, but it may help determine whether some weakly expressed archaeological features are better represented before rasterization or through alternative terrain models.

In conclusion, the thesis demonstrated that deep learning can support archaeological LiDAR interpretation, but only within a carefully controlled and expert-guided workflow. The obtained results are not sufficient for fully automatic archaeological mapping. They are, however, useful for candidate detection and prioritization, especially when the target structures have clear terrain expression. The main value of the work lies in showing how methodological decisions influence the final segmentation output and in defining the conditions under which model predictions can be meaningfully used in archaeological prospection.